\documentclass[11pt,a4paper]{article}

% ── Packages ──────────────────────────────────────────────────────────────────
\usepackage{amsmath,amsthm,amssymb}
\usepackage{hyperref}
\usepackage[margin=1in]{geometry}
\usepackage{booktabs}
\usepackage{enumitem}

% ── Theorem environments ─────────────────────────────────────────────────────
\newtheorem{theorem}{Theorem}[section]
\newtheorem{proposition}[theorem]{Proposition}
\newtheorem{lemma}[theorem]{Lemma}
\newtheorem{corollary}[theorem]{Corollary}
\newtheorem{conjecture}[theorem]{Conjecture}
\theoremstyle{definition}
\newtheorem{definition}[theorem]{Definition}
\newtheorem{assumption}[theorem]{Assumption}
\newtheorem{protocol}[theorem]{Protocol}
\theoremstyle{remark}
\newtheorem{remark}[theorem]{Remark}
\newtheorem{caveat}[theorem]{Epistemic Caveat}
\newtheorem{openproblem}{Open Problem}

% ── Macros ────────────────────────────────────────────────────────────────────
\newcommand{\LLC}{\lambda}
\newcommand{\LLChat}{\hat{\lambda}}
\newcommand{\Dadd}{\Delta_{\mathrm{add}}}
\newcommand{\Ilambda}{I_{\lambda}}
\newcommand{\Comm}{\mathrm{Comm}}
\newcommand{\Leak}{\mathrm{Leak}}
\newcommand{\Regret}{\mathrm{Regret}}
\newcommand{\Interaction}{\mathrm{Interaction}}
\newcommand{\R}{\mathbb{R}}
\newcommand{\E}{\mathbb{E}}
\newcommand{\cF}{\mathcal{F}}
\newcommand{\cC}{\mathcal{C}}
\newcommand{\cW}{\mathcal{W}}
\newcommand{\cX}{\mathcal{X}}
\newcommand{\cY}{\mathcal{Y}}

% ── Title ─────────────────────────────────────────────────────────────────────
\title{%
  SLT-Guided Residual Seeds:\\
  Weakness, Evidence Geometry, and Refinement DAGs\\
  for Transformer Adaptation%
}
\author{ProtomegaTron}
\date{July 28, 2026 --- Draft v0.1}

\begin{document}
\maketitle

\begin{abstract}
We formalize the connection between Singular Learning Theory (SLT) and the
Hyperseed framework for modular transformer adaptation.  Five claims are
developed: (1)~residual seeds are local evidence factors under a product-model
assumption with an empirically testable interaction remainder; (2)~LLC
interaction information provides a proposed diagnostic taxonomy for seed coupling;
(3)~seeds participate in refinement DAGs scored by a seven-term pregate
objective; (4)~additivity deviation $\Dadd$ is a measurable associative proxy
for compositional evidence leakage (the single novel bridge claim); and (5)~the
seed-ecology pattern extends to all LLC-estimable component-decomposable
systems (restricted universality conjecture).  A fully preregistered
experimental protocol specifies the LLC estimator, multiplicity correction via
Holm--Bonferroni, synthetic controls with pilot/confirmation split, and
falsification criteria.  All claims are framed with explicit epistemic caveats:
conditionals rather than biconditionals, association rather than causation, and
confidence intervals with family-wise error rate control.
\end{abstract}

% ══════════════════════════════════════════════════════════════════════════════
\section{Introduction}\label{sec:intro}
% ══════════════════════════════════════════════════════════════════════════════

Singular Learning Theory \cite{Watanabe2009} provides a Bayesian framework for
understanding the geometry of loss landscapes in models whose parameter-to-function
map is not injective.  The \emph{local learning coefficient} (LLC, also called
the \emph{real log canonical threshold}) $\LLC$ controls the leading correction
to the Bayesian free energy:
\begin{equation}\label{eq:free-energy}
  F_n \;=\; n\,\hat{L}_n \;+\; \LLC\,\log n
       \;-\; (m-1)\,\log\log n \;+\; O(1),
\end{equation}
where $n$ is sample size, $\hat{L}_n$ the empirical loss, and $m$ the
multiplicity of the dominant singularity.

The Hyperseed framework \cite{Weakness-SL,SLT-ResLayers} proposes that
residual corrections (``seeds'') attached to a frozen base model form an
\emph{ecology}: they are promoted, demoted, split, merged, and refined through
a directed acyclic graph of structural moves.  The central question is:
\emph{what scoring criterion should govern this ecology?}

This note argues that SLT's evidence geometry---specifically, the
decomposition of the LLC across components and the deviations from
additivity---provides the correct scoring criterion.  We formalize five claims,
each with explicit assumptions and epistemic caveats.  We then specify a
preregistered experimental protocol that makes the claims falsifiable.

\paragraph{Scope.}
This note imports definitions from the Causal Fibres framework (Note~0014)
\cite{CausalFibres,Note-0014} without re-deriving them.  The novel
contribution is Claim~4: the identification of $\Dadd$ as a quantitative
associative bridge between SLT geometry and Hyperseed compositional structure.

% ══════════════════════════════════════════════════════════════════════════════
\section{Imports from Note 0014}\label{sec:imports}
% ══════════════════════════════════════════════════════════════════════════════

The following definitions are imported from the Causal Fibres framework
\cite{CausalFibres} and Note~0014 \cite{Note-0014}.  They are cited here for
reference; see the original sources for derivations.

\begin{enumerate}[label=(\roman*),leftmargin=2em]
  \item \textbf{Frozen interface} (Def.~2.1 of \cite{CausalFibres}): a fixed
    API boundary between a base model $B$ and corrective components, through
    which information flows are defined and measurable.
  \item \textbf{Matched budget} (Def.~2.2 of \cite{CausalFibres}): a
    constraint ensuring that compared components receive equivalent
    computational resources during evaluation.
  \item \textbf{Evidence record} (Def.~2.3 of \cite{CausalFibres}): the
    accumulated Bayesian evidence for a component's contribution, scored via
    log-marginal-likelihood ratios.
  \item \textbf{Counterfactual factor closure} (Def.~2.5 of
    \cite{CausalFibres}): the condition under which a factorization of the
    evidence is closed under counterfactual interventions on component
    boundaries.
  \item \textbf{H0--H6 hypothesis ladder} (Def.~2.8 of \cite{CausalFibres}):
    a hierarchy of increasingly strong hypotheses about representation quality,
    from null (H0: no better than random) to strong compositional (H6: full
    counterfactual factor closure).
\end{enumerate}

\noindent
These imports provide the scaffold within which our SLT-based claims operate.
Claims~1--5 add LLC-based scoring criteria to 0014's structural framework; they
do not redefine the framework itself.

% ══════════════════════════════════════════════════════════════════════════════
\section{Claim 1: Residual Seeds as Local Evidence Factors}\label{sec:claim1}
% ══════════════════════════════════════════════════════════════════════════════

\subsection{Definitions}

\begin{definition}[Residual Seed]\label{def:seed}
  Let $B$ be a frozen base model with parameter space $\Theta_B$.  A
  \emph{residual seed} is a parameterized correction $r_c : \cX \to \cY$ with
  parameters $\theta_c \in \Theta_c$, such that the composite model computes
  $f(x;\theta_B,\theta_c) = B(x) + r_c(x;\theta_c)$.  A \emph{factorization}
  is a collection $\cF = \{r_1, \ldots, r_k\}$ of seeds with associated
  parameter subspaces $\Theta_{c_1}, \ldots, \Theta_{c_k}$.
\end{definition}

\begin{definition}[Evidence Neighborhood]\label{def:evidence-nbhd}
  For a seed $r_i$ in context $c$ (dataset, task, or distribution), the
  \emph{evidence neighborhood} $\cW(r_i, c)$ is the connected component of the
  set $\{\theta_i \in \Theta_i : K_{c}(\theta_i) \leq \epsilon\}$ containing
  the optimal parameters, where $K_c$ is the Kullback--Leibler divergence from
  the true distribution in context~$c$.
\end{definition}

\begin{definition}[Interaction Remainder]\label{def:interaction-remainder}
  Given a factorization $\cF = \{r_1, \ldots, r_k\}$ with joint parameter
  $\theta = (\theta_1, \ldots, \theta_k)$, define the \emph{interaction
  remainder}
  \begin{equation}\label{eq:interaction-remainder}
    R(\theta) \;=\; K(\theta) - \sum_{c=1}^{k} K_c(\theta_c),
  \end{equation}
  where $K(\theta)$ is the KL divergence of the joint model and $K_c(\theta_c)$
  is the KL divergence restricted to component~$c$.
\end{definition}

\subsection{Assumptions}

\begin{assumption}[Product-Model Assumption]\label{asm:product}
  The parameter space admits a product factorization
  $\theta = (\theta_1, \ldots, \theta_k)$ such that the prior factorizes:
  $\pi(\theta) = \prod_{c=1}^{k} \pi_c(\theta_c)$.
\end{assumption}

\begin{assumption}[Asymptotic Regime]\label{asm:asymptotic}
  The sample size $n$ is sufficiently large relative to the singularity
  structure that the free-energy asymptotic \eqref{eq:free-energy} applies
  (Watanabe's main theorem \cite{Watanabe2009}).
\end{assumption}

\begin{assumption}[Dominated Interaction Hypothesis]\label{asm:dominated}
  There exists a threshold $\eta > 0$ such that $|R(\theta)| < \eta$ for all
  $\theta$ in the evidence neighborhood of the joint model.  \textbf{This is a
  testable hypothesis, not an axiom.}
\end{assumption}

\subsection{Proposition}

\begin{proposition}[LLC Additivity under Factorization]\label{prop:llc-add}
  Under Assumptions~\ref{asm:product}--\ref{asm:dominated}, the following hold:
  \begin{enumerate}[label=(\alph*)]
    \item \textbf{LLC additivity:}
      $\displaystyle\LLC_{\mathrm{joint}} \;=\;
       \sum_{c=1}^{k} \LLC_c \;+\; O(\eta)$.
    \item \textbf{Evidence factorization:}
      $\displaystyle\log Z_n(\theta) \;=\;
       \sum_{c=1}^{k} \log Z_n(\theta_c) \;+\; O(\eta\,\log n)$.
    \item \textbf{Quasi-independence:} Each seed $r_i$ has an evidence
      neighborhood $\cW(r_i, c)$ that is approximately independent of other
      seeds' evidence neighborhoods, with coupling bounded by $O(\eta)$.
  \end{enumerate}
\end{proposition}

\begin{caveat}[Limits of the Factorization]\label{cav:claim1}
  The smallness of $R$ is a \emph{testable hypothesis}, not a sufficient proof
  of meaningful factorization.  Even if $R$ is empirically small for a given
  model, this does not establish that the decomposition is coordinate-invariant
  or structure-revealing---it establishes only that the product-prior
  decomposition is a good approximation at the observed sample size.  The
  decomposition could break under distribution shift, scale change, or
  different priors.

  The implication is conditional: factorization $\Rightarrow$ LLC additivity
  (under the assumptions).  The converse does not hold in general; LLC
  additivity could arise from cancellations in the interaction remainder.
\end{caveat}

\begin{remark}[Relation to Source Claims]\label{rem:claim1-sources}
  The evidence-factorization result builds on ``If the evidence factorizes
  across modules, the global LLC is the sum of module LLCs''
  \cite{Weakness-SL}.  The promotion criterion derives from ``Do not promote a
  residual basis because it reconstructs well.  Promote it because $R$ is
  small'' \cite{SLT-ResLayers}.  The evidence-ratio to soft-accuracy bridge
  with $O(\LLC/n)$ correction follows \cite{SLT-Accuracy}.
\end{remark}

\begin{remark}[Non-Overlap with 0014]\label{rem:nonoverlap}
  Note~0014 defines frozen interfaces, teacher-gap closure, and representation
  contests.  Claim~1 adds the LLC-additivity criterion as the promotion gate
  \emph{within} 0014's frozen-interface scaffold.  It does not redefine the
  interface.
\end{remark}

% ══════════════════════════════════════════════════════════════════════════════
\section{Claim 2: LLC Interaction Information as a Proposed Diagnostic}%
\label{sec:claim2}
% ══════════════════════════════════════════════════════════════════════════════

\begin{definition}[LLC Interaction Information]\label{def:Ilambda}
  For seeds $a, b$ in a factorization $\cF$, the \emph{LLC interaction
  information} is
  \begin{equation}\label{eq:Ilambda}
    \Ilambda(a; b) \;=\; \LLC_a + \LLC_b - \LLC_{ab},
  \end{equation}
  where $\LLC_a$, $\LLC_b$ are the per-seed LLCs and $\LLC_{ab}$ is the LLC
  of the joint $(a,b)$ subsystem.
\end{definition}

\begin{definition}[Surrogate Estimators]\label{def:surrogates}
  Two computationally tractable surrogates for $\Ilambda$:
  \begin{enumerate}[label=(\alph*)]
    \item \textbf{Mixed-Hessian norm:}
      $\|H_{ab}\| = \|\nabla_{\theta_a}\nabla_{\theta_b} L(\theta)\|_F$,
      evaluated at the MAP or posterior mean.
    \item \textbf{Commutator dominance:}
      $\Comm(a, b) = \|[J_a, J_b]\|_F$, where $J_a, J_b$ are the Jacobians
      of the component maps $r_a, r_b$ with respect to their parameters.
  \end{enumerate}
\end{definition}

\paragraph{Proposed Diagnostic Taxonomy.}
The following interpretation of $\Ilambda$ is a \emph{proposed diagnostic
framework}, not a proven classifier:

\begin{center}
\begin{tabular}{@{}lll@{}}
  \toprule
  \textbf{Sign of $\Ilambda$} & \textbf{Proposed Interpretation} & \textbf{Status} \\
  \midrule
  $\Ilambda > 0$ & Redundant/shared singular structure (merge candidate) & Hypothesis \\
  $\Ilambda < 0$ & Synergistic/emergent structure (assembly needed)      & Hypothesis \\
  $\Ilambda \approx 0$ & Factorizable independence & Consistent with Claim~1 \\
  \bottomrule
\end{tabular}
\end{center}

\begin{caveat}[Surrogates Are Not Interchangeable with $\Ilambda$]%
\label{cav:surrogates}
  The surrogates $\|H_{ab}\|$ and $\Comm(a,b)$ correlate with the true LLC
  interaction information under regularity conditions, but they measure
  different geometric quantities: local curvature cross-terms vs.\
  singularity-structure overlap.  Their agreement with $\Ilambda$ degrades near
  degenerate singularities.  Discrepancies between surrogates should be
  reported as evidence of estimator limitations, not resolved by fiat
  \cite{Weakness-SL,SLT-ResLayers}.
\end{caveat}

\begin{definition}[Surrogate Agreement]\label{def:surrogate-agreement}
  Define the surrogate agreement score
  \begin{equation}
    \delta(\Comm, \|H_{ab}\|, \Ilambda)
    \;=\; \max\!\bigl(|\Comm - \Ilambda|,\; \bigl|\|H_{ab}\| - \Ilambda\bigr|\bigr),
  \end{equation}
  after appropriate normalization.  Large $\delta$ indicates that the
  surrogates are unreliable for the pair $(a,b)$.
\end{definition}

\paragraph{Falsifiable Prediction.}
In a small transformer with known modular structure (e.g., a two-task mixture),
per-seed LLC estimates should yield $\Ilambda \approx 0$ between task-specific
seeds and $\Ilambda > 0$ between capacity-sharing seeds \cite{SLT-SubRep}.
Additionally, the surrogate measures should agree with direct LLC estimates to
within preregistered tolerances.  Surrogate--direct discrepancies, if found,
would weaken the diagnostic taxonomy rather than the underlying SLT theory.

% ══════════════════════════════════════════════════════════════════════════════
\section{Claim 3: Refinement DAGs as Operational Structure}%
\label{sec:claim3}
% ══════════════════════════════════════════════════════════════════════════════

Seeds do not exist in a fixed decomposition.  They participate in a
\emph{refinement DAG} whose vertices are factorizations and whose edges are
structural moves.

\begin{definition}[Refinement Move]\label{def:refinement-move}
  A \emph{refinement move} $\mu : \cF \to \cF'$ transforms a factorization
  $\cF$ into a successor $\cF'$ via one of the following operations:
  \textsc{split} (one seed $\to$ two), \textsc{merge} (two seeds $\to$ one),
  \textsc{mediate} (insert an intermediary), \textsc{specialize} (restrict a
  seed's domain), \textsc{demote} (remove a seed), or \textsc{transfer}
  (reassign a seed to a different context).
\end{definition}

\begin{definition}[Pregate Score]\label{def:pregate}
  The \emph{pregate score} of a factorization $\cF$ in context $c$ is the
  seven-term objective \cite{SLT-ResLayers}:
  \begin{equation}\label{eq:pregate}
    J(\cF) \;=\; n\,\hat{L}_\cF
            \;+\; \LLChat_\cF \log n
            \;-\; \log\pi(\cF)
            \;+\; \alpha\,\Comm(\cF)
            \;+\; \beta\,\Leak(\cF)
            \;+\; \gamma\,\Regret(\cF)
            \;+\; \delta\,\Interaction(\cF),
  \end{equation}
  where $\hat{L}_\cF$ is the empirical loss under factorization $\cF$,
  $\LLChat_\cF$ the estimated LLC, $\pi(\cF)$ the structural prior over
  factorizations, and the remaining terms penalize commutator dominance,
  evidence leakage, counterfactual regret, and interaction remainder,
  respectively.  The coefficients $\alpha, \beta, \gamma, \delta \geq 0$ are
  hyperparameters.
\end{definition}

\begin{definition}[Associative Signature and Regime Signature]%
\label{def:signatures}
  The \emph{associative signature} of seed $r_i$ in context $c$ is the vector
  \begin{equation}
    \sigma(r_i, c) \;=\;
      \bigl(\LLC_i,\; \Comm(r_i, \cdot),\; \Leak(r_i),\;
            \Ilambda(r_i, \cdot)\bigr) \in \R^{d_\sigma}.
  \end{equation}
  The \emph{regime signature} of a factorization $\cF$ at time $t$ is
  \begin{equation}
    \Sigma(\cF, t) \;=\;
      \bigl(\sigma(r_1, c_t), \ldots, \sigma(r_k, c_t)\bigr).
  \end{equation}
  Abrupt shifts in $\Sigma$ trigger refinement moves \cite{SLT-Regime}.
\end{definition}

\paragraph{Connection to 0014.}
Note~0014's H1 representation contest (Def.~2.8 of \cite{CausalFibres})
compares arm classes at matched budget.  Claim~3 extends this to a
\emph{dynamic} contest: the contest is not a one-shot comparison but a DAG of
refinement moves, each subject to the same evidence gates 0014 defines.

% ══════════════════════════════════════════════════════════════════════════════
\section{Claim 4: Additivity Deviation as Evidence-Leakage Proxy}%
\label{sec:claim4}
% ══════════════════════════════════════════════════════════════════════════════

\textbf{This is the single novel bridge claim.}  SLT's RLCT decomposition
across components predicts that deviations from LLC additivity signal geometric
interactions between singular regions \cite{InitSynth}.  In Hyperseed terms,
these deviations are \emph{associative signatures} of evidence leakage between
refinement stages.

\begin{definition}[Additivity Deviation]\label{def:Dadd}
  For a factorization $\cF = \{c_1, \ldots, c_k\}$, the \emph{additivity
  deviation} is
  \begin{equation}\label{eq:Dadd}
    \Dadd(\cF) \;=\;
      \Bigl|\LLC_{\mathrm{joint}}
            - \sum_{c=1}^{k} \LLC_c\Bigr|.
  \end{equation}
\end{definition}

\subsection{Epistemic Downgrades}

\begin{caveat}[No Biconditional]\label{cav:no-bicon}
  $\Dadd > 0$ is \emph{necessary but not sufficient} evidence that
  factorization $\cF$ fails to decompose the evidence geometry.  Finite-sample
  bias in LLC estimation, prior misspecification, and MCMC mixing artifacts can
  all produce non-zero $\Dadd$ even when the true singularity structure
  decomposes.

  The claim is: \emph{$\Dadd$ significantly exceeding zero, after accounting
  for estimator uncertainty, is evidence against clean factorization.}  The
  converse ($\Dadd \approx 0$ implies good factorization) is better supported
  but still conditional on the product-prior assumption
  (Assumption~\ref{asm:product}).
\end{caveat}

\begin{caveat}[Association, Not Causation]\label{cav:assoc}
  The prediction is that $\Dadd$ \emph{correlates with} graph-theoretic
  cross-stage dependency in the refinement DAG.  This is an associative claim.
  Upgrading to causal language requires an intervention that varies cross-stage
  coupling while \textbf{preserving parameter count and component
  decomposition}.  Three candidate interventions, ordered by preference:
  \begin{enumerate}[label=(\alph*)]
    \item \textbf{Coupling-strength interpolation.}  Parameterize cross-stage
      connections with a scalar $\alpha \in [0,1]$.  At $\alpha=1$ the
      architecture is unchanged; at $\alpha=0$ cross-stage information flow is
      severed.  All parameters remain in the model (zeroed weights are frozen,
      not removed).  Predict: $\Dadd$ decreases monotonically as
      $\alpha \to 0$.
    \item \textbf{Calibrated noise injection} (secondary/optional).  Add
      i.i.d.\ Gaussian noise $\varepsilon \sim \mathcal{N}(0, \sigma^2)$ at
      inter-stage boundaries, sweeping $\sigma$.  \textbf{Matching
      requirements:} at each $\sigma$, marginal activation variance and
      predictive loss must match the $\alpha$-sweep baseline within 5\%.
      Without these calibrations, noise injection is a confound.
    \item \textbf{Matched-complexity control} (if architectural changes are
      unavoidable).  Pair each structural ablation with a control: the same
      architecture but re-initialized (random) weights, matching parameter
      count and decomposition.
  \end{enumerate}
  Until an intervention is run, ``$\Dadd$ tracks cross-stage dependency'' means
  association.
\end{caveat}

\subsection{Confidence-Interval Requirement}

All reported $\Dadd$ values must include:
\begin{itemize}
  \item Bootstrap or MCMC-derived confidence intervals on individual
    $\LLC_c$ estimates.
  \item Propagated uncertainty on $\Dadd$ itself (via the delta method or
    bootstrap on the sum).
  \item Multiplicity-corrected significance: $\Dadd$ is declared ``meaningfully
    non-zero'' iff its Holm-adjusted $p$-value $< 0.05$ across the full
    comparison family (see \S\ref{subsec:multiplicity}).
\end{itemize}

\subsection{Falsifiable Experimental Check}

Take a small transformer (or synthetic model with known singular structure).
Compute per-layer LLC estimates via the \texttt{devinterp} local learning
coefficient estimator with stated hyperparameters.  Report $\Dadd$ with
confidence intervals.  Independently compute a graph-theoretic cross-stage
dependency score on the refinement DAG.  Test for statistically significant
correlation (Spearman $\rho$, Holm-adjusted $p < 0.05$).

\begin{itemize}
  \item If $\Dadd$ is significantly elevated but the DAG shows clean stage
    separation, the associative claim is \textbf{falsified}.
  \item If they correlate, the Hyperseed--SLT bridge has non-trivial empirical
    support---but \emph{not} causal evidence unless an intervention experiment
    (Caveat~\ref{cav:assoc}) is also run.
\end{itemize}

\paragraph{Why Novel.}
The individual pieces (LLC decomposition \cite{Weakness-SL,SLT-SubRep},
refinement DAGs \cite{Weakness-SL}, counterfactual factor closure
\cite{CausalFibres}) exist in the source papers.  The novel claim is that
$\Dadd$ is the \emph{quantitative observable} that connects SLT geometry to
Hyperseed compositional structure---as a measurable, falsifiable associative
bridge.

% ══════════════════════════════════════════════════════════════════════════════
\section{Claim 5: Restricted Adaptation Conjecture}\label{sec:claim5}
% ══════════════════════════════════════════════════════════════════════════════

The seed-ecology pattern---base process + controlled corrective factors, scored
by evidence geometry, refined through a DAG of structural moves, monitored for
regime change---extends beyond residual layers to a defined class of systems.

\begin{definition}[LLC-Estimable Component-Decomposable System]%
\label{def:lcd-system}
  A system belongs to the \emph{LLC-estimable component-decomposable} class iff:
  \begin{enumerate}[label=(C\arabic*)]
    \item\label{cond:param} It admits a finite parametric description
      $\theta \in \Theta \subset \R^d$.
    \item\label{cond:analytic} Its loss function $L(\theta)$ is real-analytic
      (or admits a resolution of singularities).
    \item\label{cond:llc} Local learning coefficients are estimable at each
      component (the posterior concentrates sufficiently for MCMC-based LLC
      estimation to converge).
    \item\label{cond:decomp} It has an explicit component decomposition
      $\cF = \{c_1, \ldots, c_k\}$ with defined parameter subspaces
      $\Theta_{c_i}$ and a product-prior factorization.
    \item\label{cond:R} The interaction remainder $R(\theta)$ from
      Definition~\ref{def:interaction-remainder} is well-defined and estimable.
  \end{enumerate}
\end{definition}

\begin{conjecture}[Restricted Hyperseed Universality]\label{conj:universality}
  Any system satisfying conditions \ref{cond:param}--\ref{cond:R} that
  maintains:
  \begin{enumerate}[label=(\roman*)]
    \item a base predictive process,
    \item an ecology of modular corrective seeds scored by local evidence
      (LLC-additivity checks, not reconstruction error alone), and
    \item a refinement DAG over candidate factorizations
  \end{enumerate}
  will exhibit the same LLC-additivity / interaction-information /
  regime-signature dynamics described in Claims~1--4.
\end{conjecture}

\paragraph{Conjectured Properties within this Class.}
\begin{itemize}
  \item \textbf{Incremental compression:} corrective components are compression
    features for the base model's error stream.
  \item \textbf{Evidence gating:} promotion/demotion is governed by
    LLC-additivity checks, not reconstruction error alone.
  \item \textbf{Regime detection:} LLC/signature shifts trigger structural
    revision \cite{SLT-Regime}.
  \item \textbf{Interaction diagnostics:} coupling type is diagnosed via
    $\Ilambda$ (Claim~2's proposed taxonomy).
  \item \textbf{Goal-stability limits:} SLT diagnostics detect structural
    shifts but not semantic truth; the semantics route via distinctions
    $\to$ partitions $\to$ symmetries $\to$ singularities adds the missing
    layer \cite{SLT-GoalStab,SLT-Semantics}.
\end{itemize}

\begin{caveat}[Scope Exclusions]\label{cav:scope}
  Explicitly excluded from Conjecture~\ref{conj:universality}: non-parametric
  models, models where LLC estimation does not converge, systems without a
  natural component decomposition, and infinite-dimensional parameter spaces
  (unless a finite-dimensional effective description exists).
\end{caveat}

% ══════════════════════════════════════════════════════════════════════════════
\section{Preregistered Experimental Protocol}\label{sec:protocol}
% ══════════════════════════════════════════════════════════════════════════════

Before running the falsification experiment for Claims~1, 2, and~4, the
following must be preregistered.

% ─── 8.1 ─────────────────────────────────────────────────────────────────────
\subsection{LLC Estimator Specification}\label{subsec:estimator}

\begin{protocol}[LLC Estimator]\label{prot:llc}
  \leavevmode
  \begin{itemize}
    \item \textbf{Estimator:} \texttt{devinterp} local learning coefficient
      estimator (Lau et al., 2024).
    \item \textbf{MCMC method:} Stochastic Gradient Langevin Dynamics (SGLD)
      with specified step-size schedule, number of chains, burn-in length, and
      thinning interval.  All hyperparameters are fixed before data collection.
    \item \textbf{Convergence diagnostic:} $\hat{R} < 1.05$ across chains;
      effective sample size $\geq 200$ per component.
    \item \textbf{Hyperparameters to freeze:} learning rate, temperature
      schedule, number of posterior samples, number of chains, burn-in
      duration, and thinning interval.
  \end{itemize}
\end{protocol}

% ─── 8.2 ─────────────────────────────────────────────────────────────────────
\subsection{Uncertainty Quantification and Multiplicity Correction}%
\label{subsec:multiplicity}

\paragraph{Per-component uncertainty.}
\begin{itemize}
  \item \textbf{Per-component $\LLC_c$:} Bootstrap confidence intervals
    ($B \geq 1000$ resamples) or MCMC posterior credible intervals (95\%).
  \item \textbf{$\Dadd$:} Propagated via the delta method or bootstrap on the
    sum; reported with 95\% CI.
  \item \textbf{$\Ilambda(a, b)$:} Same uncertainty propagation; reported
    with 95\% CI.
  \item \textbf{Unadjusted significance threshold:} $\Dadd$ is declared
    ``meaningfully non-zero'' iff the 95\% CI excludes zero.
\end{itemize}

\paragraph{Comparison family.}
The experiment involves simultaneous statistical tests.  Define the comparison
family as:
\begin{itemize}
  \item \textbf{$\Dadd$ tests:} $k$ coupling levels $\times$ $m$ components
    $= k \cdot m$ comparisons.
  \item \textbf{$\Ilambda$ tests:} $\binom{m}{2} = m(m-1)/2$ pairwise
    interaction-information comparisons.
  \item \textbf{Total family size:}
    \begin{equation}\label{eq:family-size}
      N \;=\; k \cdot m \;+\; \frac{m(m-1)}{2}.
    \end{equation}
    With $k = 11$ ($\alpha$ from 0.0 to 1.0 in steps of 0.1) and $m$
    components (fixed before data collection), $N$ is stated explicitly in
    the preregistration.
\end{itemize}

\paragraph{Holm--Bonferroni Procedure.}
The Holm--Bonferroni sequential rejection procedure (Holm, 1979) controls the
family-wise error rate (FWER) at $\alpha_{\mathrm{family}} = 0.05$:
\begin{enumerate}
  \item Order all $N$ $p$-values:
    $p_{(1)} \leq p_{(2)} \leq \cdots \leq p_{(N)}$.
  \item Reject $H_{0,(j)}$ iff
    $p_{(j)} \leq \alpha_{\mathrm{family}} / (N - j + 1)$ for all $j' \leq j$.
  \item Equivalently, report Holm-adjusted $p$-values:
    \begin{equation}\label{eq:holm}
      \tilde{p}_{(j)} \;=\;
        \max_{j' \leq j}\bigl\{(N - j' + 1)\cdot p_{(j')}\bigr\}.
    \end{equation}
  \item An individual comparison is declared significant iff its Holm-adjusted
    $p$-value $< 0.05$.
\end{enumerate}

\paragraph{Adjusted Confidence Intervals.}
Where CIs are reported alongside Holm-corrected tests, use Bonferroni-adjusted
CIs at level $1 - \alpha_{\mathrm{family}}/N$ per comparison (conservative but
consistent with FWER control).  These replace unadjusted 95\% CIs in all
multi-comparison contexts.

\paragraph{Rationale for Holm over Alternatives.}
Holm--Bonferroni controls FWER (appropriate here because a single false
positive about coupling structure would undermine the bridge claim) while being
strictly more powerful than Bonferroni.  FDR-controlling methods
(Benjamini--Hochberg) would be appropriate for screening many components for
follow-up, but here each comparison bears directly on the claim.

% ─── 8.3 ─────────────────────────────────────────────────────────────────────
\subsection{Component-to-DAG Mapping}\label{subsec:dag-mapping}

\begin{protocol}[Component-to-DAG Mapping]\label{prot:dag}
  \leavevmode
  \begin{itemize}
    \item \textbf{Component definition:} Each residual block (or attention
      head / MLP sublayer, depending on granularity---to be fixed before
      experiment) is one component $c_i$.
    \item \textbf{DAG construction:} Edges from $c_i$ to $c_j$ iff
      information flow exists (measured by gradient attribution or activation
      patching).  Edge weights from mutual information or ablation impact.
    \item \textbf{Cross-stage dependency score:} Weighted edge density between
      non-adjacent DAG stages.  The formal graph-theoretic definition is stated
      before data collection.
  \end{itemize}
\end{protocol}

% ─── 8.4 ─────────────────────────────────────────────────────────────────────
\subsection{Synthetic Controls with Pilot/Confirmation Split}%
\label{subsec:controls}

\paragraph{Split Design.}
Synthetic control data is generated once and split into two disjoint sets
before any analysis:
\begin{itemize}
  \item \textbf{Pilot calibration set (50\%):} Used \emph{exactly once} for
    threshold calibration (\S\ref{subsec:thresholds}).  After thresholds are
    frozen, this data is permanently set aside.
  \item \textbf{Held-out confirmation set (50\%):} Reserved exclusively for
    admission tests.  Never examined during calibration.
  \item \textbf{Contingency:} If the pilot set is too small for stable
    threshold estimation (bootstrap CV of threshold estimates $> 0.2$), the
    split shifts to 60/40 pilot/confirmation.  This contingency is
    preregistered.
\end{itemize}

\paragraph{Null/Negative Control.}
Two independent MLPs sharing no parameters, trained on disjoint tasks from the
same data distribution.  Components are a priori independent by construction.
\begin{itemize}
  \item \textbf{Expected:} $\Dadd \approx 0$, $\Ilambda \approx 0$ between
    the two networks, $R < \eta$.
  \item \textbf{Gate:} If the estimator reports non-zero $\Dadd$ on this null
    control, the estimator or its hyperparameters are miscalibrated; do not
    proceed.
\end{itemize}

\paragraph{Positive Control.}
Two MLPs (component$_1$, component$_2$) connected by a fixed linear coupling
layer $W_{\mathrm{mix}}$ of known rank~$r$, trained jointly on the same task
mixture.  The component decomposition is defined a priori:
component$_1 = \text{MLP}_1$ parameters, component$_2 = \text{MLP}_2$
parameters, coupling $= W_{\mathrm{mix}}$.

\begin{caveat}[No Analytic Proportionality]\label{cav:no-prop}
  The prediction is: \emph{$\Dadd$ is expected to be significantly nonzero
  (Holm-adjusted $p < 0.05$) on the positive control, with the effect size
  estimated from the pilot calibration set.}  No analytic proportionality
  between $\Dadd$ and
  $\mathrm{rank}(W_{\mathrm{mix}}) \cdot \|W_{\mathrm{mix}}\|_2$
  is claimed.  Computing the actual RLCT contribution
  of the coupling requires a resolution of singularities for the specific loss
  surface, which has not been performed.  If a tighter analytic bound becomes
  available (e.g., via explicit resolution of the $W_{\mathrm{mix}}$ singularity
  under a Gaussian prior), it will be stated as a separate, independently
  testable prediction.
\end{caveat}

\paragraph{Decision Rule (Confirmation Set Only).}
Both controls must pass before claims about the test model are made:
\begin{itemize}
  \item Null: Holm-adjusted $p > 0.05$ for $\Dadd$.
  \item Positive: Holm-adjusted $p < 0.05$ for $\Dadd$.
\end{itemize}

% ─── 8.5 ─────────────────────────────────────────────────────────────────────
\subsection{Model and Data}\label{subsec:model}

\begin{itemize}
  \item \textbf{Test model:} Small transformer ($\leq 6$ layers) trained on a
    two-task mixture.
  \item \textbf{Held-out model:} A second, architecturally distinct model
    (e.g., a 4-layer transformer if the test model is 6 layers, or a
    convolutional residual network) trained on the same data.  Predictions for
    this held-out model are registered before fitting.  If Claims~1--4 hold on
    the test model but fail on the held-out model, the claims are
    model-specific, not general.
  \item \textbf{Training data:} Fixed dataset, split specified before training.
  \item \textbf{LLC estimation data:} Training split (for loss-landscape
    exploration).
  \item \textbf{Evaluation data:} Held-out split, not used for LLC estimation
    or estimator tuning.
\end{itemize}

% ─── 8.6 ─────────────────────────────────────────────────────────────────────
\subsection{Numerical Decision Thresholds}\label{subsec:thresholds}

All significance tests across the $\alpha$-coupling sweep and component
comparisons are corrected for multiplicity via the Holm--Bonferroni procedure
(\S\ref{subsec:multiplicity}).  Thresholds are calibrated on the pilot set
(\S\ref{subsec:controls}) and frozen before confirmation data is examined.

\begin{itemize}
  \item \textbf{$\Dadd$ significance:} $\Dadd$ is declared ``meaningfully
    non-zero'' iff its Holm-adjusted $p$-value $< 0.05$.  Reported with
    Bonferroni-adjusted CI at level $1 - 0.05/N$.
  \item \textbf{Interaction remainder threshold:}
    $\eta = 0.1 \times \LLC_{\max}$, where $\LLC_{\max}$ is the largest
    per-component LLC.  $R/\LLC_{\max} < 0.1$ counts as ``dominated.''
  \item \textbf{$\Ilambda$ significance:} $|\Ilambda(a,b)|$ is declared
    non-negligible iff its Holm-adjusted $p$-value $< 0.05$.
  \item \textbf{Correlation threshold (Claim~4 primary test):} Spearman
    $\rho$ with Holm-adjusted $p < 0.05$ (one-tailed).  The minimum detectable
    $\rho$ is set using the pilot calibration set, targeting $\geq 80\%$ power
    at the pilot-observed effect size.  If the pilot effect is too small for
    80\% power, report achieved power and note the study is underpowered.
  \item \textbf{Sign accuracy (Claim~2):} Agreement between
    $\mathrm{sign}(\Ilambda)$ and known modular structure.  Threshold set at
    $2\sigma$ above null expectation (random assignment baseline), calibrated
    on pilot data.  Cohen's $\kappa$ threshold frozen before confirmation.
  \item \textbf{Factorization prevalence (Claim~1):} $R < \eta$ for a
    preregistered fraction of components.  Threshold set at the ROC-optimal
    discrimination point from pilot data.  If controls show perfect separation,
    use 80\% (conservative).
\end{itemize}

\paragraph{Circularity Prevention.}
The three calibration steps each consume the pilot set.  Once thresholds are
frozen, the pilot data is permanently set aside.  All subsequent
control-pass/fail decisions and claim evaluations use only the held-out
confirmation set.  No threshold is revised after examining confirmation data.

% ─── 8.7 ─────────────────────────────────────────────────────────────────────
\subsection{Analysis Plan}\label{subsec:analysis}

\begin{itemize}
  \item \textbf{Primary test (Claim~4):} Spearman correlation between
    per-component $\Dadd$ and cross-stage dependency score.  Report
    coefficient, $p$-value, and CI.
  \item \textbf{Secondary test (Claim~2):} Agreement between
    $\mathrm{sign}(\Ilambda)$ and known modular structure.  Report accuracy,
    Cohen's $\kappa$.
  \item \textbf{Tertiary test (Claim~1):} Fraction of components where
    $R < \eta$ ($\eta$ preregistered).  Report proportion and bootstrap CI.
\end{itemize}

\paragraph{Falsification Criteria.}
\begin{itemize}
  \item Claim~4 is falsified if the Holm-adjusted $p$-value for the
    $\Dadd$--dependency correlation exceeds 0.05.
  \item Claim~2's taxonomy is weakened if sign accuracy falls below the
    pilot-calibrated threshold.
  \item Claim~1's factorization hypothesis is weakened if the proportion of
    components with $R < \eta$ falls below the pilot-calibrated prevalence
    threshold.
  \item All thresholds are those frozen after pilot calibration
    (\S\ref{subsec:thresholds}); none are adjusted post hoc.
\end{itemize}

\paragraph{Intervention Design (Eventual Causal Upgrade of Claim~4).}
\begin{itemize}
  \item \textbf{Primary:} Coupling-strength interpolation ($\alpha$ sweep from
    1.0 to 0.0 in steps of 0.1) at selected stage boundaries, preserving
    parameter count and decomposition.  Re-estimate $\Dadd$ at each setting;
    test for monotonic decrease.  Holm--Bonferroni applied across all $\alpha$
    levels and components jointly.
  \item \textbf{Secondary (optional):} Calibrated noise injection ($\sigma$
    sweep) with mandatory marginal-activation-variance and predictive-loss
    matching (Caveat~\ref{cav:assoc}(b)).  If architectural ablation is used,
    include a matched-complexity control.
  \item This is planned as a follow-up, not part of the initial
    preregistration.
\end{itemize}

% ══════════════════════════════════════════════════════════════════════════════
\section{Open Questions}\label{sec:open}
% ══════════════════════════════════════════════════════════════════════════════

\begin{openproblem}[Coordinate Invariance]
  The product-prior decomposition (Assumption~\ref{asm:product}) is
  coordinate-dependent.  Under what conditions does the LLC-additivity result
  (Proposition~\ref{prop:llc-add}) hold in reparameterized coordinates?  Is
  there a coordinate-free characterization of ``good factorization''?
\end{openproblem}

\begin{openproblem}[Resolution of Singularities for Coupled Components]
  The positive control (\S\ref{subsec:controls}) uses a coupling matrix
  $W_{\mathrm{mix}}$.  What is the RLCT of the joint model as a function of
  $\mathrm{rank}(W_{\mathrm{mix}})$ and the component RLCTs?  An explicit
  resolution would convert Claim~4's associative prediction into a quantitative
  one.
\end{openproblem}

\begin{openproblem}[Scaling of the Interaction Remainder]
  How does $R(\theta)$ scale with model depth, width, and dataset size?  If $R$
  grows faster than $\LLC_{\max}$, the ``dominated interaction'' hypothesis
  (Assumption~\ref{asm:dominated}) fails at scale, and the factorization
  framework may not apply to large models.
\end{openproblem}

\begin{openproblem}[Semantic Registration]
  Claim~5 notes that SLT diagnostics detect structural shifts but not semantic
  truth \cite{SLT-GoalStab}.  How should the distinctions $\to$ partitions
  $\to$ symmetries $\to$ singularities pipeline \cite{SLT-Semantics} be
  formalized to provide the missing semantic layer?
\end{openproblem}

\begin{openproblem}[Beyond Real-Analytic Loss]
  Condition~\ref{cond:analytic} of Definition~\ref{def:lcd-system} restricts
  to real-analytic loss functions.  Many practical losses (e.g., those involving
  ReLU activations) are only piecewise analytic.  Can the framework be extended
  via stratified resolution of singularities, or does the LLC decomposition
  fundamentally require analyticity?
\end{openproblem}

\begin{openproblem}[Dynamic Pregate Coefficients]
  The pregate score \eqref{eq:pregate} has fixed coefficients
  $\alpha, \beta, \gamma, \delta$.  Should these be adapted over the
  refinement DAG (e.g., via meta-learning on the DAG trajectory)?  What are
  the convergence implications of a non-stationary scoring function?
\end{openproblem}

% ══════════════════════════════════════════════════════════════════════════════
% Bibliography
% ══════════════════════════════════════════════════════════════════════════════
\begin{thebibliography}{99}

\bibitem{Weakness-SL}
B.~Goertzel.
\newblock Weakness as local evidence: LLC decomposition, evidence
  factorization, quantale message-passing, and refinement DAGs.
\newblock Working paper, 2025--2026.

\bibitem{SLT-ResLayers}
B.~Goertzel.
\newblock {SLT} and residual layers: Seven-term pregate formula, computable
  surrogates, mixed-{H}essian and commutator diagnostics.
\newblock Working paper, 2025--2026.

\bibitem{SLT-Accuracy}
B.~Goertzel.
\newblock {SLT}-accuracy-weakness v1: Evidence-ratio to soft-accuracy bridge
  with $O(\lambda/n)$ correction.
\newblock Working paper, 2025--2026.

\bibitem{SLT-SubRep}
B.~Goertzel.
\newblock {SLT} for {SubRep} v5: Interaction information interpretation and
  interaction complexity.
\newblock Working paper, 2025--2026.

\bibitem{SLT-Regime}
B.~Goertzel.
\newblock {SLT} for regime-change detection: Module-wise {LLC} signature shifts
  as refinement triggers.
\newblock Working paper, 2025--2026.

\bibitem{SLT-GoalStab}
B.~Goertzel.
\newblock {SLT}-goal-stability v4: Semantic truth limits of {SLT} diagnostics.
\newblock Working paper, 2025--2026.

\bibitem{SLT-Semantics}
B.~Goertzel.
\newblock {SLT}-semantics v2: Distinctions, partitions, symmetries, and
  singularities.
\newblock Working paper, 2025--2026.

\bibitem{SLT-Evolution}
B.~Goertzel.
\newblock {SLT}-evolution: Free-energy geodesics and estimation of distribution
  algorithms.
\newblock Working paper, 2025--2026.

\bibitem{CausalFibres}
B.~Goertzel.
\newblock Causal fibres 0.4.0: {H0--H6} hypothesis ladder, frozen interfaces,
  and representation contests.
\newblock Technical report, 2025--2026.

\bibitem{Note-0014}
B.~Goertzel.
\newblock Note 0014: Emotion regime operators.
\newblock Hyperseed working notes, 2026.

\bibitem{InitSynth}
B.~Goertzel.
\newblock Initial synthesis: Deviations from additivity signal geometric
  interactions between singular regions.
\newblock Working paper, 2025--2026.

\bibitem{Watanabe2009}
S.~Watanabe.
\newblock \emph{Algebraic Geometry and Statistical Learning Theory}.
\newblock Cambridge University Press, 2009.

\bibitem{Holm1979}
S.~Holm.
\newblock A simple sequentially rejective multiple test procedure.
\newblock \emph{Scandinavian Journal of Statistics}, 6(2):65--70, 1979.

\bibitem{Lau2024}
E.~Lau, Z.~Murfet, and others.
\newblock \texttt{devinterp}: Local learning coefficient estimation via {SGLD}.
\newblock Software package, 2024.

\end{thebibliography}

\end{document}
